\section{How the Church Learned to Keep Christmas}\label{sec:history}

The twelve days of the 1962 Christmas liturgy did not drop from heaven as a
finished book. They grew across sixteen centuries of Roman worship, imperial
politics, patristic preaching, and papal reform. Knowing when each layer entered
the Missal and Breviary turns puzzling details---why a Greek martyr is named in
the Dawn Mass, why the Pope moved across Rome before sunrise, why 1~January
keeps both Mary and the Circumcision, and why 1962 vests the Holy Innocents in
red---into intelligible history.

\subsection{25 December in Rome and the East}

The oldest surviving documentary proof of Christmas on 25~December stands at the
very head of the Roman \textit{Depositio Martyrum} inside the Chronograph of
354, reflecting Roman liturgical usage already fixed by at least~336:
\textit{VIII kal.\ Ian.\ natus Christus in Betleem Iudeae}---``Eighth day before
the Kalends of January [25~December]: Christ was born in Bethlehem of Judea.''
Before the fourth century, as C.~C.~Martindale notes in his 1908 historical
survey, writers such as Clement of Alexandria recorded varied speculative dates
in spring for Christ's birth, while Eastern churches celebrated the Nativity
together with the Baptism and the Magi on 6~January under the title of Epiphany
or Theophany.

From Rome, the 25~December feast spread eastward during the late fourth and
fifth centuries. On 25~December~380, St.~Gregory of Nazianzus preached
Oration~38 in Constantinople, calling the feast both ``Theophany'' and
``Birthday'' because God was manifested to man by birth, while promising his
hearers that ``a little later on you will see Jesus submitting to be purified in
the River Jordan''---proof that Constantinople already kept 25~December apart
from 6~January. Within a decade, Antioch received the Roman date, where
St.~John Chrysostom preached around~386--388 that the feast had been known there
for less than ten years; Jerusalem still kept the Nativity on 6~January when the
pilgrim Egeria visited around~385, adopting 25~December in the fifth century,
and Egypt followed around~432.

Historians have long weighed two principal explanations for Rome's choice of
25~December. One points to Emperor Aurelian's 274 solstice festival, the
\textit{Dies Natalis Solis Invicti} (``Birthday of the Unconquered Sun''), which
the Church countered by proclaiming Christ the true Sun of Justice. The other,
argued by Louis Duchesne, traces 25~December to an early Christian calculation
that placed the Annunciation and the Passion alike on 25~March, yielding a birth
nine months later on 25~December. Neither explanation is settled dogma; both may
have reinforced one another.

What fifth-century primary evidence does prove beyond doubt is that pagan solar
customs still clung to Roman converts on Christmas morning. Preaching in
Sermon~27 (section~4) around the middle of the fifth century, Pope St.~Leo the
Great rebuked a habit he had witnessed with his own eyes at the Vatican:

\begin{quote}
even some Christians think it is so proper to do this that, before entering the
blessed Apostle Peter's basilica, which is dedicated to the One Living and true
God, when they have mounted the steps which lead to the raised platform, they
turn round and bow themselves towards the rising sun and with bent neck do
homage to its brilliant orb. We are full of grief and vexation that this should
happen, which is partly due to the fault of ignorance and partly to the spirit
of heathenism.
\end{quote}

Even if some claimed to honor the Creator through his bright creature, Leo
refused to let Christian worship borrow the gestures of pagan solstice rites.
The true Light had been born in Bethlehem, not risen in the winter sky.

\subsection{The Three Masses and Their Roman Stations}

By the fifth and sixth centuries, Roman prayer-books---first the Roman
\textit{libelli} preserved in the Veronese Sacramentary, then the Old Gelasian
Sacramentary and the Gregorian \textit{Hadrianum}---had gathered a rich cycle of
Nativity collects and secrets. By the end of the sixth century, the Roman bishop
celebrated three distinct Masses on Christmas Day. Pope St.~Gregory the Great
supplies our earliest patristic proof in the opening sentence of Homily~8 on the
Gospels (c.~590--591), preached at St.~Mary Major: ``By God's mercy we are to
say three Masses to-day, so that there is not much time left for preaching.''

The Gregorian \textit{Hadrianum} preserves the three ancient papal stations
that shaped the entire Western rite:

\begin{enumerate}
\item \textbf{First Mass at Midnight:} \textit{Ad Sanctam Mariam maiorem ad
Praesepe}, at the Basilica of St.~Mary Major ``at the Crib.'' Rebuilt by Pope
St.~Sixtus~III (432--440) immediately after the Council of Ephesus proclaimed
Mary Mother of God in~431, St.~Mary Major became Rome's own Bethlehem, housing
its chapel of the Crib in imitation of Jerusalem's nocturnal vigil at Bethlehem.
Bl.~Ildefonso Schuster recounts in \textit{The Sacramentary} that during
Midnight Mass in that very chapel on Christmas~1075, the violent Roman noble
Cencio burst in and dragged Pope St.~Gregory~VII wounded from the altar---only
for the Roman people to storm Cencio's tower at dawn and carry the pope back to
St.~Mary Major to finish the Christmas liturgy.
\item \textbf{Second Mass at Dawn:} \textit{Ad Sanctam Anastasiam}, at the
Basilica of St.~Anastasia at the foot of the Palatine Hill. Why did the pope
descend from the Esquiline at daybreak to a church dedicated to a fourth-century
martyr of Sirmium? As Schuster's Roman stational archaeology explains,
St.~Anastasia's stood beside the Palatine residences of the Byzantine imperial
officials and garrison in Rome; her relics had been translated to Constantinople
and her feast fell on 25~December. The pope therefore went down at dawn to
celebrate the stational Mass for the Byzantine court on its patroness's day. As
imperial rule in Rome faded, the liturgy was absorbed into Christmas as the
\textit{Missa in aurora}---leaving behind the unique Commemoration of
St.~Anastasia still printed in the 1962 Dawn Mass.
\item \textbf{Third Mass in the Day:} Originally \textit{Ad Sanctum Petrum},
the great papal stational Mass at St.~Peter's Basilica on the Vatican, where
Pope St.~Leo~III crowned Charlemagne on Christmas Day~800. By the later Middle
Ages and in the Tridentine and 1962 Missals, the Day Mass station shifted back
to St.~Mary Major, while the Epiphany station remained at St.~Peter's.
\end{enumerate}

St.~Francis of Assisi popularized the living crib at Greccio in December~1223
(a turning point in parish and household devotion), yet Roman churches had built
liturgical cribs centuries earlier---not only at St.~Mary Major, but also under
Pope Gregory~IV (827--844) at Santa Maria in Trastevere.

\subsection{The Octave, 1 January, and the Most Holy Name}

Both the Old Gelasian Sacramentary and the Gregorian \textit{Hadrianum} already
set St.~Stephen (26~December), St.~John the Evangelist (27~December), and the
Holy Innocents (28~December) immediately after 25~December, assigning stations
at St.~Stephen on the Caelian Hill, St.~Mary Major, and St.~Paul Outside the
Walls. Medieval liturgists---Sicard of Cremona in his \textit{Mitrale}, John
Beleth in his \textit{Summa}, and William Durandus in Book~VII, Chapter~42 of
his \textit{Rationale}---expounded these three days as the \textit{Comites
Christi}, the newborn King's closest companions, followed in the thirteenth
century by St.~Thomas of Canterbury (29~December) beside Pope St.~Sylvester~I
(31~December).

The history of 1~January forms a four-layer palimpsest still visible in the
1962 Missal. First, the fourth- through sixth-century Church kept 1~January as a
day of protest, fasting, and prayer against the riotous pagan New Year Kalends:
the \textit{Roman Martyrology} honors St.~Almachius slain by gladiators in Rome
for crying out against the idol sacrifices on the Octave, St.~Peter Chrysologus
preached \textit{Sermon~155} against the Kalends processions, the Old Gelasian
Sacramentary set a penitential Mass \textit{Prohibendum ab Idolis} directly
beside \textit{In Octavas Domini}, and the Second Council of Tours (567)
prescribed a three-day fast with a \textit{Missa Circumcisionis}. Second, before
Rome received the Byzantine Marian feasts in the seventh century, 1~January was
Rome's oldest Marian feast---\textit{Natale Sanctae Mariae}, celebrated at
Santa Maria in Trastevere (and in the \textit{Hadrianum} at the Pantheon,
\textit{ad S.~Mariam ad martyres}), which left 1~January its Marian Collect
(\textit{Deus, qui salutis aeternae}) and its \textit{O admirabile commercium}
antiphons. Third, Gallican and medieval usage fused that Marian Octave with the
Feast of the Circumcision of the Lord; John Beleth records in the twelfth
century that Rome and many churches even offered two Masses on 1~January, one of
the Circumcision and one of Our Lady.

The Feast of the Most Holy Name of Jesus, rooted in St.~Bernard's twelfth-century
preaching on the Song of Songs and widely propagated in the late Middle Ages,
stood on the Second Sunday after the Epiphany in pre-1913 Roman books (such as
the 1862 Ratisbon Missal and the 1888 Pustet Breviary). On 23~October~1913, Pope
St.~Pius~X's motu proprio \textit{Abhinc duos annos} and the accompanying Sacred
Congregation of Rites decree of 28~October~1913 permanently transferred the Holy
Name into Christmastide---to the Sunday falling between 1~January and 6~January,
or to 2~January whenever no Sunday intervenes---anchoring it beside the naming
of the Child on the eighth day.

\subsection{From the 1955 Simplification to the 1962 Books}

The 1962 \textit{Missale Romanum} and \textit{Breviarium Romanum} are governed
by Pope St.~John~XXIII's motu proprio \textit{Rubricarum instructum}
(25~July~1960) and the \textit{Codex Rubricarum} (26~July~1960, effective
1~January~1961), building on Pope Pius~XII's 1955 simplification. Readers
accustomed to older hand missals notice six clear changes in the 1962 Christmas
season:

\begin{enumerate}
\item \textbf{One octave only:} Under General Rubric~64, only the octaves of
Christmas (second class), Easter (first class), and Pentecost (first class)
remain. The pre-1955 simple octaves of St.~Stephen (2~January), St.~John
(3~January), and the Holy Innocents (4~January), together with the Vigil of the
Epiphany (5~January) and the Epiphany octave, were abolished.
\item \textbf{Sunday within the Octave wins over 26--28 December:} Under the
1960 table of precedence (RG~91, row~15 vs.\ row~16), when 26, 27, or
28~December falls on a Sunday, the Sunday within the Octave is celebrated with
an ordinary commemoration of St.~Stephen, St.~John, or the Holy Innocents in Low
Mass only, instead of being displaced by those feasts to 30~December as in
pre-1960 books.
\item \textbf{Holy Innocents vested in red with joyful chants:} Pre-1960 books
kept 28~December on weekdays in mourning violet without the \textit{Gloria},
\textit{Te Deum}, or \textit{Alleluia}, using folded chasubles and
\textit{Benedicamus Domino}. In 1962, folded chasubles are abolished (RG~137),
red is worn unconditionally (RG~124~e), and the \textit{Gloria}, \textit{Te
Deum}, \textit{Alleluia} (\textit{Laudate, pueri}), and \textit{Ite, missa est}
(RGMR~507) are said whenever the feast is celebrated.
\item \textbf{29 and 31 December restored to the Octave:} St.~Thomas of
Canterbury (29~December) and St.~Sylvester~I (31~December) were reduced to
commemorations (\textit{Variationes} n.~5), leaving 29, 30, and 31~December
uniformly as Days~5, 6, and~7 within the Octave (using the hybrid Mass
\textit{Puer natus est} with the Dawn readings).
\item \textbf{1 January restored as Octave Day, First Class:} Under
\textit{Variationes} nn.~6 and~12, 1~January was renamed \textit{In Octava
Nativitatis Domini} and raised to the first class (RG~91, row~5), while keeping
the Circumcision Gospel (Luke~2:21) and its ancient Marian prayers intact.
\item \textbf{Simplified Sunday Matins:} On the Sunday within the Octave,
Breviary Rubrics~162~a and~220 reduce Matins to one nocturn of nine psalms and
three lessons (Romans~1 and St.~Ambrose on Luke~2:33), while Christmas Day,
St.~Stephen, St.~John, the Holy Innocents, 1~January, and the Most Holy Name
retain three nocturns and nine lessons.
\end{enumerate}

The postconciliar calendar approved by Pope St.~Paul~VI in \textit{Mysterii
Paschalis} (14~February~1969) later titled 1~January the Solemnity of Mary,
Mother of God, and extended Christmas Time as one season through the Baptism of
the Lord. In the 1962 Roman Rite, by contrast, \textit{Tempus Nativitatis} keeps
its crisp twelve-day arc from First Vespers of 24~December through None of
5~January, followed by \textit{Tempus Epiphaniae} from First Vespers of
5~January through 13~January---a continuous organic growth across the centuries
of the kind Pope Benedict~XVI described in his 22~December~2005 address to the
Roman Curia as reform in the continuity of the one Church.
